% Modernised reading — fonds Alexandre Grothendieck, shelfmark 29, batch 6
% (pages 101-120).
%
% This is an INTERPRETATION, not a transcription. It re-reads the pages in
% today's notation and vocabulary, states the hypotheses the manuscript leaves
% implicit, and drops the critical apparatus so that the mathematics reads
% straight through. Everything the brackets and underlines carried moves into a
% footnote. In any dispute of substance the transcription governs: whoever
% cites Grothendieck must cite batch-06.fr.tex.
%
% Scope: the folder was read WHOLE before any of this was written — all eleven
% batches, pages 1-216, in one pass — and one file was then written per batch.
% Read at that scope, this batch is the folder's toolbox: page 103 is the
% monodromy theorem whose relation f g f^{-1} = g^q the tame quotient of a
% local Galois group satisfies; pages 112-115 write out the Bertini lemma that
% batch 4 USED at page 66 without stating; pages 116-122 turn the warning of
% the 1967 letter — « be careful that the inertia groups are determined only up
% to interior automorphism! » (page 7) — into a theorem.
%
% Departures from the page, each footnoted where it occurs:
%   - Lemme 6.5 of page 103 says « endomorphismes » while using F^{-1} and
%     placing the eigenvalues in the units: F and G must be invertible;
%   - Remarque 4 of page 112 writes pi_1^!(X_t) -> pi_1^!(X_t) on both sides
%     of an arrow whose left member must be Y_t;
%   - the inertia class of page 116 depends on the chosen extension of the
%     covering, which the page says and then fixes by convention; the
%     convention is restated here where the class is first used.
%
% Pass: Opus 5 (claude-opus-5), 2026-08-29 — first pass, unchecked against the
% pages by a human.
\documentclass[11pt,a4paper]{article}
% Le préambule commun à toutes les transcriptions du fonds Grothendieck.
%
% Il tient en une idée : **l'incertitude se compose, elle ne se résout pas.**
% Une transcription qui lisse un mot douteux a détruit la seule chose pour
% laquelle on la faisait — un lecteur doit pouvoir savoir, à chaque endroit,
% ce qui était sur le feuillet et ce qui a été ajouté. Les six macros qui
% suivent sont donc l'appareil critique, et non de la décoration.
%
% Les mêmes commandes sont reconnues par `scripts/render.mjs`, qui produit la
% vue de lecture du site. Une macro ajoutée ici doit l'être aussi là-bas,
% faute de quoi le rendu échouera bruyamment — ce qui est voulu : un rendu
% silencieusement faux est pire qu'un rendu absent.

\ProvidesPackage{grothendieck}[2026/08/08 Transcriptions du fonds Grothendieck]

\usepackage{amsmath,amssymb,amsthm}
\usepackage{mathrsfs}
\usepackage{stmaryrd}
% Les diagrammes commutatifs sont partout dans ce fonds : c'est la langue
% même de la géométrie algébrique de Grothendieck, et une transcription qui
% les rendrait en prose ne transcrirait rien.
\usepackage{tikz-cd}
% Français par défaut — c'est la langue des feuillets — mais l'anglais est
% chargé aussi : la traduction et le résumé sont des documents anglais, et la
% typographie française (espaces avant les deux-points) y serait une faute.
% Ces documents basculent par \selectlanguage{english} en tête de corps.
\usepackage[english,french]{babel}
\usepackage{fontspec}
\usepackage{geometry}
\geometry{a4paper,margin=25mm,marginparwidth=28mm}
\usepackage{marginnote}
\usepackage{xcolor}
% ulem, not soul. `soul`'s \st and \ul are built on a character-by-character
% scan that XeLaTeX's font machinery breaks: the document fails with an
% undefined control sequence rather than degrading. ulem does the same two jobs
% and is Unicode-engine safe. `normalem` stops it hijacking \emph, which the
% transcriptions use for Grothendieck's own emphasis.
\usepackage[normalem]{ulem}
\usepackage{hyperref}
\hypersetup{colorlinks=true,linkcolor=black,urlcolor=black,pdfborder={0 0 0}}

% --- Métadonnées du lot -----------------------------------------------------
% Reprises telles quelles par le script de rendu, qui les affiche en tête de
% la vue de lecture. Elles fixent ce que le fichier prétend couvrir : sans
% elles, une transcription flotte sans point d'attache dans un fonds de seize
% mille feuillets.

\newcommand{\folder}[1]{\gdef\grothFolder{#1}}
\newcommand{\batch}[1]{\gdef\grothBatch{#1}}
\newcommand{\leaves}[2]{\gdef\grothFirst{#1}\gdef\grothLast{#2}}
\newcommand{\foldertitle}[1]{\gdef\grothFoldertitle{#1}}
\gdef\grothFolder{?}\gdef\grothBatch{?}\gdef\grothFirst{?}\gdef\grothLast{?}\gdef\grothFoldertitle{}

% --- Le résumé --------------------------------------------------------------
%
% Il ouvre la lecture modernisée, sous le titre « Résumé », et il est composé
% en petit. Ce n'est pas un ornement : c'est le seul passage du document qui
% ne soit pas des mathématiques, et un lecteur doit pouvoir le reconnaître
% avant de l'avoir lu — puis le sauter s'il connaît déjà le sujet, ou s'y
% arrêter s'il ne le connaît pas. Le filet qui le suit marque la reprise du
% corps.

\newenvironment{resume}
  {\small\setlength{\parskip}{0.45em}}
  {\par\medskip\hrule\medskip}

% --- Mots-clés ---------------------------------------------------------------
%
% En anglais, en clôture du résumé de la lecture modernisée : le vocabulaire
% moderne sous lequel chercher ce que les feuillets construisent. Le manifeste
% du site (`npm run manifest`) les extrait pour étiqueter la cote entière —
% cette ligne est la source unique des tags, il n'y a pas de fichier à côté.
\newcommand{\keywords}[1]{\par\smallskip\noindent
  {\footnotesize\textsc{Keywords} --- \textit{#1}}\par}

% --- L'appareil critique ----------------------------------------------------

% Le numéro de feuillet, en marge. C'est l'ancre qui permet de régler un
% désaccord sur une formule en regardant *un* feuillet plutôt qu'une chemise
% de six cents. Le site s'en sert aussi pour tourner les pages du fac-similé
% quand on fait défiler la transcription.
\newcommand{\leaf}[1]{%
  \marginnote{\footnotesize\color{gray!70}\textsf{#1}}%
  \label{leaf:#1}\ignorespaces}

% Illisible : on ne devine pas. Un mot inventé qui se lit comme les autres est
% le pire résultat possible.
\newcommand{\ill}{\textcolor{red!60!black}{[\,\ldots\,]}}

% Lecture douteuse : le mot est donné, et signalé comme incertain.
\newcommand{\uncertain}[1]{\uline{#1}}

% Ajout de l'éditeur — une lettre manquante, un mot sous-entendu.
\newcommand{\add}[1]{\textcolor{blue!50!black}{[#1]}}

% Biffé par Grothendieck lui-même. Ce qu'il a rayé fait partie du document :
% une rature dit souvent où la pensée a changé de direction.
\newcommand{\struck}[1]{\sout{#1}}

% Note du transcripteur, sur l'état matériel du feuillet.
\newcommand{\note}[1]{\par\smallskip{\small\color{gray!60!black}\textit{[#1]}}\par\smallskip}

% Note marginale de Grothendieck, distincte du corps du texte.
\newcommand{\marginal}[1]{\par\smallskip\noindent\hspace{1em}%
  {\small\color{gray!60!black}#1}\par\smallskip}

% --- Titre ------------------------------------------------------------------

\AtBeginDocument{%
  \begin{center}
    {\large\bfseries Fonds Alexandre Grothendieck --- cote n\textsuperscript{o}~\grothFolder}\\[2pt]
    {\itshape\grothFoldertitle}\\[4pt]
    {\small feuillets \grothFirst--\grothLast\ (lot \grothBatch)}\\[2pt]
    {\footnotesize\color{gray!60!black}Transcription établie d'après le fac-similé numérisé
     par l'Université de Montpellier.\\
     Les lectures incertaines sont soulignées, les illisibles notées [\,\ldots\,],
     les ajouts entre crochets.}
  \end{center}
  \vspace{1em}}

\endinput


\folder{29}
\batch{6}
\pages{101}{120}
\dating{1959-1969}
\watermark{Édition de démonstration\\interprétation personnelle de l'œuvre}
\foldertitle{Groupe fondamental [Autour de SGA 4 et SGA 7] — lecture modernisée}

\begin{document}

\section*{Résumé}

\begin{resume}
Ce lot est la boîte à outils du carton. Il n'a pas de thèse ; il a des
instruments, et chacun sert ailleurs.

Le premier est le plus célèbre, et il tient en dix lignes. Considérez un objet
géométrique qui dépend d'un paramètre et qui se casse quand le paramètre
atteint zéro. Deux mouvements l'agitent : la monodromie, qui est ce que fait
l'objet quand on tourne autour de zéro, et le Frobenius, qui est l'ombre portée
de l'arithmétique quand le corps de base est fini. Ces deux mouvements ne
commutent pas : le Frobenius \emph{multiplie} le tour par le nombre
d'éléments du corps résiduel. C'est la relation $f g f^{-1} = g^{q}$, et
Grothendieck en tire, en une demi-page, que la monodromie ne peut pas être
quelconque : ses valeurs propres sont des racines de l'unité. Un objet
algébrique ne peut pas tourner indéfiniment sans revenir sur lui-même. C'est
le \emph{théorème de monodromie locale}, et la page qui le porte est barrée de
deux traits diagonaux, comme si elle avait été recopiée ailleurs.

Le deuxième instrument est le lemme de Bertini, écrit ici pour de bon --- le
lot 4 s'en servait sans l'énoncer. Une variété se coupe par un espace linéaire
générique, et la section reste irréductible, unibranche, normale, lisse, selon
qu'on veille sur le lieu où les choses tournent mal. Le corollaire est la
surjectivité des groupes fondamentaux, et le théorème 6 en donne enfin la
version \emph{modérée} : pour un $L_{t}$ qui n'est plus générique, mais où le
diviseur induit reste à croisements normaux. La remarque 7 est la seule du
carton où Grothendieck dise franchement où il bute : en caractéristique 2 et en
dimension impaire, il ne sait pas démontrer la conjugaison des cycles
évanescents pour un pinceau qui ne soit pas générique, parce que la
représentation devient sauvage.

Le troisième est le plus utile et le moins glorieux : la \emph{fonctorialité
des groupes d'inertie}. En 1967, page 7 de ce carton, Grothendieck avertissait
Murre : « be careful that the inertia groups are determined only up to interior
automorphism ! », et le théorème 6.4 de Murre s'était effondré là-dessus. Ces
quatre pages font de l'avertissement un énoncé : le groupe d'inertie n'est pas
un sous-groupe, c'est une classe de conjugaison ; il ne dépend pas seulement du
point mais du prolongement choisi du revêtement ; sous une hypothèse de
régularité il est quotient d'un produit de modules de Tate, un facteur par
composante du diviseur ; et un morphisme le transforme selon la matrice
d'entiers qui exprime l'image inverse des diviseurs. La condition pour que
cette matrice soit inversible est jolie : il faut que son déterminant soit,
au signe près, une puissance de la caractéristique.

Entre ces trois, le lot contient quatre pages isolées --- un symbole local sur
une variété abélienne, une redite du théorème kummérien du lot 5, une
définition de $\pi_{2}$, et de l'algèbre commutative sur les polynômes unitaires
--- et un feuillet du tapuscrit d'EGA IV, corrigé de sa main et portant en
marge le meilleur avertissement de la boîte : « Attention à l'ambiguïté dans
\emph{lisse} ! ».

Les noms modernes à chercher sont \emph{théorème de monodromie locale},
\emph{quasi-unipotence}, \emph{groupe d'inertie}, \emph{pinceau de Lefschetz},
\emph{cycles évanescents}.

\keywords{local monodromy theorem, quasi-unipotence, inertia group, Lefschetz pencil, vanishing cycles}
\end{resume}

\pagerange{101}{102}
\section*{Deux pages isolées : un symbole local, et le théorème kummérien redit (pages 101 et 102)}

\subsection*{Le symbole local}

La page 101 est le carbone d'un tapuscrit français, sans lien avec le reste du
carton. Elle achève la démonstration d'une réciprocité pour les symboles locaux
sur une variété abélienne $A$ : pour deux cycles $\underline{a}$,
$\underline{b}$ et un diviseur $X$, on veut
\[
  j_{o}(X_{\underline{a}}, \underline{b}) \;=\;
  j_{o}(X^{-}_{\underline{b}}, \underline{a}) .
\]
Le cas où $\underline{b}$ appartient au noyau d'Albanese se traite par
linéarité et par l'invariance du symbole sous les translations, ce qui ramène
au cas déjà traité ; le cas général se ramène, en modifiant $\underline{b}$, à
celui où les supports sont étrangers, et l'on conclut par l'invariance du
symbole $i$ sous l'application $\zeta(x) = a + b - x$.

C'est de la théorie des symboles locaux au sens de Serre, et le carton n'y
revient pas.

\subsection*{Le théorème kummérien, troisième énoncé}

La page 102 --- numérotée 3 par lui, dans un cercle --- reprend en corollaire le
programme de la page 100. Elle est écrite très vite et raturée d'un bout à
l'autre ; la transcription en laisse une large part illisible et cette lecture
n'invente pas ce qui manque. Ce qui se lit est l'articulation, et c'est elle
qui compte.

Soit $X$ un schéma normal localement intègre, $(D_{i})$ des diviseurs de
Cartier, et $X'$ un revêtement dont les groupes d'inertie relatifs aux $D_{i}$
sont tels que $X'$ soit modérément ramifié le long des $D_{i}$. On a alors une
extension
\[
  1 \longrightarrow I \longrightarrow G \longrightarrow \Gamma
  \longrightarrow 1 ,
\]
avec $\tilde{X} = X'/I$ revêtement principal de $X$ de groupe $\Gamma$, d'où
une classe $\alpha \in H^{2}(\Gamma, I)$, puis une classe
\[
  \alpha \in H^{2}\bigl(X,\ \tilde{X} \times^{\Gamma} I\bigr) .
\]
On dispose d'un isomorphisme canonique
$\tilde{D} \times^{\Gamma} I \simeq \prod_{i} \mu_{n_{i},\,D}$ sur $D$,
compatible aux opérations de $\Gamma$, et la conclusion visée est la même que
celle des pages 89, 100 et 200 : la classe est celle du fibré normal,
\[
  c\bigl(\mathcal{O}_{X}(D)\bigr) \in H^{2}(X, \mu_{\underline{n}}) ,
  \qquad
  c\bigl(T_{D/X}\bigr) \in H^{2}(D, \mu_{\underline{n}}) .
\]

\pagerange{103}{104}
\section*{Le théorème de monodromie locale (pages 103 et 104)}

\subsection*{La relation, et d'où elle vient}

La page 103 est barrée de deux longs traits diagonaux, et c'est la seule page
du lot qui se lise proprement de bout en bout.

Le contexte est celui d'un corps local de corps résiduel fini à $q$ éléments.
Le quotient modéré de son groupe de Galois est engendré par deux éléments : un
relèvement $f$ du Frobenius, et un générateur $g$ de l'inertie modérée --- qui,
après choix d'isomorphismes $\mathbb{Z}_{\ell}(1) \simeq \mathbb{Z}_{\ell}$
pour tout $\ell$, engendre topologiquement
$\prod_{\ell \neq p} \mathbb{Z}_{\ell}$. Ils satisfont la relation
\[
  (6.4) \qquad f\,g\,f^{-1} \;=\; g^{q} ,
\]
qui exprime que le Frobenius multiplie le tour par $q$. Le groupe est le
produit semi-direct de $\mathbb{Z}$ par $\prod_{\ell\neq p}\mathbb{Z}_{\ell}$,
$\mathbb{Z}$ opérant par multiplication par $q$.

\subsection*{Le lemme}

\textbf{Lemme.} Soient $F$, $G$ deux \emph{automorphismes} d'un espace
vectoriel $E$ de dimension finie sur un corps, et $q > 1$ un entier tels que
\[
  F G F^{-1} \;=\; G^{q} .
\]
Alors $G$ est \emph{quasi-unipotent} : ses valeurs propres sont toutes des
racines de l'unité ; autrement dit, il existe $N \geqslant 1$ tel que $G^{N}$
soit unipotent.\footnote{La page écrit « deux endomorphismes », tout en
utilisant $F^{-1}$ dans la relation et en plaçant l'ensemble $\Phi$ des valeurs
propres dans le groupe multiplicatif $\bar{R}^{*}$. Il faut $F$ inversible pour
que la relation ait un sens, et $G$ inversible pour la conclusion : si $0$
était valeur propre de $G$, il le serait de $G^{q}$ et l'argument ne
l'excluerait pas, alors que $0$ n'est pas racine de l'unité. Les deux
hypothèses sont supplées ici.}

\emph{Démonstration.} L'ensemble $\Phi$ des valeurs propres de $G$ dans une
clôture algébrique du corps de base est aussi celui de $FGF^{-1}$, donc celui
de $G^{q}$ : c'est donc une partie finie de $\bar{R}^{*}$ stable par
$\lambda \mapsto \lambda^{q}$. Comme $\Phi$ est fini, pour tout
$\lambda \in \Phi$ il existe deux entiers $N' \neq N''$ avec
$\lambda^{q^{N'}} = \lambda^{q^{N''}}$, d'où $\lambda^{q^{N'} - q^{N''}} = 1$ :
$\lambda$ est une racine de l'unité.

C'est le \emph{théorème de monodromie locale}, dans sa forme purement
algébrique : la relation de commutation entre Frobenius et inertie interdit à
la monodromie d'être quelconque. Il n'y a pas d'hypothèse de géométrie dans le
lemme --- seulement la relation $(6.4)$ et la finitude de la
dimension.\footnote{C'est le lemme qui figure dans SGA 7, exposé I, comme
première étape du théorème de monodromie ; il y est attribué à Grothendieck et
la démonstration est celle-ci. Que la page du carton soit barrée de deux
diagonales n'est pas un jugement sur l'énoncé : c'est la marque, fréquente dans
ce carton, d'un texte reporté ailleurs.}

\subsection*{$\pi_{2}$}

La page 104 définit, en cinq lignes, l'objet dont le lot 5 s'était servi sans
le définir.

\begin{enumerate}
  \item[a)] Si $X$ est simplement connexe, le foncteur
    $M \rightsquigarrow H^{2}(X, M)$, sur les groupes abéliens finis, est
    exact --- $H^{1}(X, M)$ étant nul --- donc pro-représentable, et l'on pose
    \[
      H^{2}(X, M) \;=\; \operatorname{Hom}_{\mathrm{cont}}
      \bigl(\pi_{2}(X),\, M\bigr) .
    \]
  \item[b)] Dans le cas général, $\pi_{2}(X, \xi) = \pi_{2}\bigl(\tilde{X}(\xi)
    \bigr)$, où $\tilde{X}(\xi)$ est le revêtement universel pointé défini par
    $\xi$.
\end{enumerate}

C'est le $\pi_{2}$ de l'homotopie étale, et le procédé --- définir un groupe
d'homotopie comme le pro-objet qui pro-représente un foncteur de cohomologie
--- est exactement celui que la lettre à Serre de 1959 emploie pour $\pi_{1}$
(page 126) et que les pages 164 à 185 systématiseront. Le bas de la page porte
un croquis en perspective de trois diviseurs se coupant à croisements normaux ;
il n'est pas repris, n'étant pas un diagramme.

\pagerange{105}{110}
\section*{Un feuillet d'EGA IV, et de l'algèbre commutative (pages 105 à 110)}

\subsection*{« Attention à l'ambiguïté dans \emph{lisse} ! »}

La page 105 est le feuillet IV-878 du tapuscrit d'EGA IV, portant ses
corrections. Il court de la proposition (17.4.10) --- un morphisme localement
de présentation finie est non ramifié si et seulement si ses fibres le sont ---
jusqu'au théorème (17.5.1), qui donne les quatre caractérisations équivalentes
de la lissité en un point : lissité, platitude plus lissité de la fibre,
régularité du morphisme, et lissité formelle de
$\mathcal{O}_{X,x}$ comme $\mathcal{O}_{Y,y}$-algèbre pour les topologies
discrètes.

Deux corrections sont substantielles : « localement \emph{noethérien} » est
biffé et remplacé par « localement \emph{de type fini} » dans l'argument sur la
fibre, et les deux équivalences sont étiquetées en marge. La troisième est un
avertissement, et il vaut pour tout le carton : « Attention à l'ambiguïté dans
« lisse » ! » --- le mot ayant servi, dans les premières rédactions, tantôt
pour \emph{lisse} et tantôt pour \emph{simple}.\footnote{L'ambiguïté n'est pas
théorique : la page 74 du carton, dans « Pic birationnel », écrit
« $X$ propre et \emph{simple} sur $k$ » là où le corollaire 1 demande
\emph{lisse}, et le corollaire 2 de la même page demande un « modèle propre et
\emph{simple} ». C'est l'un des rares endroits où l'avertissement de la marge
d'EGA porte sur le carton lui-même.}

\subsection*{Le $\pi_{1}$ local le long de deux branches}

La page 106 dessine deux branches $X_{T'}$ et $X_{T''}$ se coupant en un point
$p$, et un petit disque $Z = \operatorname{Spec}
\hat{\mathcal{O}}_{X,a}$ autour du croisement. Les deux traces
$Z' = X_{T'} \cap Z$ et $Z'' = X_{T''} \cap Z$ correspondent aux corps
résiduels $k(p')$ et $k(p'')$, qui sont les corps locaux en $a$ des deux
branches. Le diagramme de groupes fondamentaux qui en résulte,

\begin{tikzcd}[column sep=small, row sep=small, nodes={font=\scriptsize}]
 & \hat{\mathbb{Z}} = \pi_{1}(Z') \arrow[dl] \arrow[dr] & \pi_{1}(Z'') = \hat{\mathbb{Z}}
 \arrow[dl] \arrow[d] \\
\pi_{1}(X_{T'}) \arrow[dr] & \pi_{1}(X_{T''}) \arrow[d] & \pi_{1}(Z) =
\hat{\mathbb{Z}}^{2} \arrow[dl] \\
 & \pi_{1}(X_{T}) &
\end{tikzcd}

est celui qu'un théorème de descente devrait clore, et la page le dit sans le
faire : « On doit prouver un th.\ de descente de nature essentiellement locale,
plus déterminer certains hom.\ canoniques ». C'est la situation locale des
pages 41 et 87, vue une troisième fois.

\subsection*{Polynômes unitaires et idéaux étrangers}

Les pages 107, 109 et 110 sont de l'algèbre commutative élémentaire, mise en
place pour un usage qui n'est pas dit. Le point est le suivant. Pour deux
idéaux $Q$, $Q'$ d'un anneau commutatif $A$, on a l'équivalence
\[
  Q + Q' = A \;\Longleftrightarrow\; Q \cap Q' = Q\,Q' ,
\]
et sous cette condition $A/QQ' \to A/Q \times A/Q'$ est un isomorphisme. D'où :
les décompositions de $B/R$ en produit de deux anneaux correspondent aux
couples d'idéaux $J$, $J'$ avec $JJ' = R$ et $J + J' = B$ ; et si $B = A[T]$ et
$R = (P)$ avec $P$ unitaire, ces décompositions sont exactement les
factorisations $P = QQ'$ en polynômes unitaires. L'unicité permet de se ramener
au cas local, et l'existence se voit alors sur le complété.

C'est le mécanisme de factorisation dont on se sert pour décomposer une algèbre
finie étale en facteurs connexes ; le carton s'en servira à la page 152, où la
proposition 3.2 caractérise les morphismes finis étales de rang $n$.

\pagerange{112}{115}
\section*{« Théorème de Bertini et variantes » (pages 112 à 115)}

Un tapuscrit corrigé de sa main. Il énonce ce que la page 66 --- « $\pi_{1}$
birationnel », lot 4 --- utilisait sans le démontrer, puis va plus loin.

Dans tout ce qui suit $k$ est un corps et $\mathbf{P}^{r} = \mathbf{P}^{r}_{k}$.

\subsection*{Les deux énoncés de base}

\textbf{Théorème de Bertini 1.} Soit $f : X \to \mathbf{P}^{r}_{k}$ un
morphisme, $X$ un schéma algébrique géométriquement irréductible, et
$n = \dim \overline{f(X)}$. Soit $L_{t}$ le sous-espace linéaire
\emph{générique} de codimension $s \leqslant n-1$. Alors
$Y_{t} = f_{t}^{-1}(L_{t})$, dans $X_{t} = X \otimes_{k} k(t)$, est
géométriquement irréductible.

\textbf{Énoncé local 2.} Soit $f : X \to \mathbf{P}^{r}$ avec $X$ algébrique.
Soit $Z \subset X$ l'ensemble des points où $f$ est ramifié \emph{ou} où $X$
n'est pas géométriquement unibranche --- respectivement pas normal, pas lisse.
Si $\dim \overline{f(Z)} < s$, alors $Y_{t} = f_{t}^{-1}(L_{t})$ est
géométriquement unibranche --- respectivement normal, lisse.

Le second est le mécanisme de tout ce qui suit, et il faut voir sa forme : la
propriété qu'on veut sur la section générique s'obtient en \emph{écartant} par
la dimension le lieu où elle échoue sur $X$.

\subsection*{Les corollaires}

\textbf{Corollaire 3.} Sous les conditions de 1 et 2, avec $k$ algébriquement
clos, le morphisme canonique $\pi_{1}(Y_{\bar{t}}) \to \pi_{1}(X)$ est
\emph{surjectif}.

\textbf{Remarque 4.} Utilisant les théorèmes de changement de base pour le
$\pi_{1}$, on conclut que
\[
  (\ast) \qquad \pi_{1}(Y_{\bar{t}}) \longrightarrow \pi_{1}(X_{\bar{t}})
\]
est surjectif \emph{si $X$ est propre}, et que
\[
  (\ast\ast) \qquad \pi_{1}^{!}(Y_{\bar{t}}) \longrightarrow
  \pi_{1}^{!}(X_{\bar{t}})
\]
est surjectif sans hypothèse de propreté, où $!$ désigne le plus grand quotient
profini premier à l'exposant caractéristique de $k$.\footnote{Le tapuscrit
écrit, pour $(\ast\ast)$, « $\pi_{1}^{!}(X_{\bar{t}}) \to
\pi_{1}^{!}(X_{\bar{t}})$ », les deux membres identiques. Le membre de gauche
doit être $Y_{\bar{t}}$, comme en $(\ast)$ ; c'est une faute de frappe et elle
est corrigée ici.} On \emph{n'affirme pas} que $(\ast)$ soit toujours surjectif
sans propreté : c'est exactement la différence que $(\ast\ast)$ mesure, et elle
est du côté sauvage.

\textbf{Corollaire 4 (de 1).} Si $X$ est de plus propre, alors pour
\emph{toute} sous-variété linéaire $L_{t}$ de codimension $s$ ---
générique ou non --- $Y_{t} = f_{t}^{-1}(L_{t})$ est géométriquement connexe.
C'est le théorème de connexité.

\textbf{Corollaire 5.} Sous les hypothèses de 3, avec $X$ propre,
$\pi_{1}(Y_{\bar{t}}) \to \pi_{1}(X)$ est surjectif, l'identification
$\pi_{1}(X) \simeq \pi_{1}(X_{\bar{t}})$ venant du théorème de changement de
base. En marge : « cf.\ SGA 1 X 2.11 ».

\subsection*{Le théorème 6 : la version modérée, pour $t$ non générique}

C'est le résultat propre du tapuscrit, et il répond au besoin qu'énonce la page
113 : un énoncé de surjectivité pour un $t$ \emph{non} générique, sans
hypothèse de propreté, et pour des quotients plus gros que ceux de
$(\ast\ast)$.

\textbf{Théorème 6.} Supposons $k$ algébriquement clos, $X$ propre,
$f : X \to \mathbf{P}^{r}$, $L_{t}$ un sous-espace linéaire de codimension $s$
défini sur $k$, et $D \subset X$ un sous-schéma fermé, tels que :

\begin{enumerate}
  \item[a)] $Y_{t} = f^{-1}(L_{t})$ est non vide ;
  \item[b)] pour tout $x \in Y_{t}$ : $X$ est lisse en $x$, $Y_{t}$ est lisse
    en $x$ et de codimension $s$ dans $X$ en $x$, et $f$ est net en $x$ ;
  \item[c)] pour tout $x \in D_{t} = D \cap Y_{t}$ : ou bien
    $\operatorname{codim}_{x}(D_{t}, Y_{t}) \geqslant 2$, ou bien $D$ est un
    diviseur sur $X$ en $x$ qui induit sur $Y_{t}$ en $x$ un diviseur à
    croisements normaux.
\end{enumerate}

Supposons $U = X - D$ irréductible et posons $U_{t} = U \cap Y_{t}$. Soit
$\pi^{t}_{1}(U)$ le quotient de $\pi_{1}(U)$ classifiant les revêtements
modérément ramifiés le long des composantes de codimension 1 de $D$ qui
rencontrent $Y_{t}$, et $\pi^{t}_{1}(U_{t})$ le quotient analogue pour
$D_{t}$. Alors
\[
  \pi^{t}_{1}(U_{t}) \longrightarrow \pi^{t}_{1}(U)
\]
est \emph{surjectif}.

\emph{Démonstration.} Pour $L_{t'}$ générique la conclusion est le corollaire
3 : un revêtement étale connexe $V$ de $U$ reste connexe au-dessus de
$U_{t'}$. Il faut la faire descendre à $t$. Sur un voisinage ouvert $T$ de
$\bar{t}$ dans la grassmannienne, la famille des $Y_{t'}$ est un morphisme
propre et lisse $Y \to T$, et celle des $D_{t'}$ un sous-schéma fermé $\Delta$
de $Y$, réunion d'un $\Delta_{1}$ de codimension $\geqslant 2$ sur chaque
fibre --- qui ne compte pas, par le théorème de pureté relative --- et d'un
diviseur $\Delta_{2}$ à croisements normaux relatifs. Le revêtement $V$
définit un revêtement $V'$ de $U' = Y - \Delta$ ; par pureté, $V'$ se prolonge
en un revêtement étale $V''$ de $U'' = Y - \Delta_{2}$, modérément ramifié sur
les fibres relativement à $\Delta_{2}$. On conclut par le \emph{théorème de
spécialisation pour le groupe fondamental modéré}.\footnote{Le tapuscrit
renvoie à « SGA 1 XIII (exposé de Mme Raynaud) », au futur : « qui
figure(ra) ». C'est le second renvoi du carton à ce théorème, après celui de la
page 22 ; les deux le citent comme l'outil qui manque, et la ramification
modérée au sens large de la page 21 est ce qui le rend applicable.}

\subsection*{Remarque 7 : où il bute}

Avec les notations de SGA 7 XVIII, le théorème 6 s'applique à tout pinceau de
Lefschetz $D$ et donne la surjectivité de
\[
  \pi^{t}_{1}(D - \check{D}) \longrightarrow
  \pi^{t}_{1}(\mathbf{P} - \check{X}) ,
\]
d'où le théorème de conjugaison des cycles évanescents --- donc
l'irréductibilité de la cohomologie évanescente --- \emph{pourvu qu'on ne soit
pas dans le cas $\operatorname{car} k = 2$ avec $\dim X$ impaire}.

Dans ce cas-là, dit-il, « je ne sais pas prouver la conjugaison sous
$\pi_{1}(D - \check{D})$, sauf pour un pinceau \emph{générique}, puisqu'alors
la représentation envisagée est sauvage ; j'ignore si elle reste vraie pour un
pinceau \emph{général} ». Il ajoute que la conjugaison des sous-groupes de
monodromie locale \emph{dans} $\operatorname{Aut}(H^{n-1}(Y))$, qui est très
élémentaire, reste vraie et suffit pour le théorème caractérisant le cas où la
condition (A) n'est pas vérifiée.

C'est l'un des deux endroits du carton où Grothendieck écrive noir sur blanc
qu'il ne sait pas ; l'autre est la page 179.

\pagerange{116}{120}
\section*{« Comportement fonctoriel des groupes d'inertie » (pages 116 à 120)}

Un second tapuscrit corrigé, en cinq numéros dont le cinquième déborde sur le
lot 7. Il traite ce que la lettre de 1967 signalait comme la cause de
l'effondrement du théorème 6.4 de Murre.

\subsection*{1. Ce qu'est un groupe d'inertie, et de quoi il dépend}

Soient $X$ un schéma, $U$ un ouvert, $G$ un groupe profini, $U'$ un revêtement
principal de $U$ de groupe $G$, $x$ un point géométrique de $X$, $\bar{X}$ le
localisé strict de $X$ en $x$, $\bar{U}$ l'image inverse de $U$ et
$\bar{U}' = U' \times_{U} \bar{U}$. Donnons-nous de plus un $\bar{X}$-schéma
$\bar{X}'$ à opérateurs $G$ prolongeant $\bar{U}'$.

Le cas d'intérêt est celui où $x \in X - U$ est adhérent à $U$, où $\bar{X}'$
est entier sur $\bar{X}$, et où
$\mathcal{O}_{\bar{X}} = \bar{p}_{*}(\mathcal{O}_{\bar{X}'})^{G}$ --- de sorte
que $G$ opère transitivement sur la fibre $\bar{X}'_{\bar{x}}$. Si $X$ est
normal en $x$ et $\bar{U}$ non vide, on prend pour $\bar{X}'$ le
\emph{normalisé} de $\bar{X}$ relativement à $\bar{U}'$, comme limite projective
des normalisés dans les quotients finis.

Pour $\xi \in \bar{X}'_{\bar{x}}$, le stabilisateur
\[
  G_{\xi} \;=\; \{\, g \in G \;:\; g\cdot\xi = \xi \,\}
\]
est le \emph{groupe d'inertie relatif à $\xi$}. Quand $G$ opère transitivement
sur la fibre, les $G_{\xi}$ forment une classe de sous-groupes conjugués, qu'on
appelle par abus \emph{les} sous-groupes d'inertie relatifs à $x$.

Il faut lire l'avertissement qui suit, parce qu'il est le sujet de tout ce
tapuscrit : \emph{cette classe n'est pas déterminée par la seule connaissance
de $x$}, elle dépend du choix du prolongement $\bar{X}'$ de $\bar{U}'$. La
convention --- $X$ normal en $x$, $x$ spécialisation d'un point de $U$,
prolongement par normalisation --- est ce qui la fixe, et elle est sous-entendue
partout ensuite.

\subsection*{2. Fonctorialité}

Soit une seconde situation $(Y, V, H, V')$, un morphisme $f : X \to Y$ avec
$f(U) \subset V$, un morphisme $f' : U' \to V'$ au-dessus de lui, et un
homomorphisme $\varphi : G \to H$ rendant $f'$ équivariant. Le cas courant est
celui où $U'$ et $V'$ sont des quotients des revêtements universels de $U$ et
$V$ pointés en $a$ et $b = f(a)$, où $G$ et $H$ sont des quotients de
$\pi_{1}(U,a)$ et $\pi_{1}(V,b)$, et où $\varphi$ est induit par
$f_{*}$ ; alors $f'$ et $\varphi$ sont déterminés sans ambiguïté par $f$, et
leur existence signifie seulement que $f_{*}$ passe aux quotients.

Étant donné un prolongement $\bar{f}' : \bar{X}' \to \bar{Y}'$, on a alors la
relation évidente
\[
  (2.1) \qquad \varphi(G_{\xi}) \;\subset\; H_{\bar{f}'(\xi)} :
\]
un morphisme envoie inertie dans inertie. Et si les deux actions sont
transitives sur les fibres, $\varphi$ applique tout groupe d'inertie de $G$
relatif à $x$ dans un groupe d'inertie de $H$ relatif à $y = f(x)$.

Reste l'existence de $\bar{f}'$. En pratique $\bar{U}'$ est schématiquement
dense dans $\bar{X}'$ et $\bar{Y}'$ est séparé sur $\bar{Y}$, ce qui rend
$\bar{f}'$ unique s'il existe ; et il existe --- en l'interprétant comme
$\bar{X}' \to \bar{Y}' \times_{\bar{Y}} \bar{X}'$ --- lorsque $\bar{Y}'$ est
entier sur $\bar{Y}$, par EGA II 6.1.14. La commutation aux actions de $G$ est
alors automatique.

\subsection*{3. Le calcul local}

Supposons $\mathcal{O}_{X,x}$ régulier et $\bar{U} = \bar{X} - \bar{D}$ avec
$\bar{D}$ un diviseur à croisements normaux dans le schéma strictement local
régulier $\bar{X}$, et supposons la ramification de $\bar{U}'$ \emph{modérée}
aux points maximaux de $\bar{D}$. Le calcul local classique donne alors : pour
tout $\xi$, le groupe d'inertie $G_{\xi}$ est canoniquement un \emph{quotient}
du groupe
\[
  (3.1) \qquad I \;=\; \prod_{\ell \neq p}
  \underline{Z}_{\ell}(1)\bigl(k(\bar{x})\bigr)^{C}
  \;=\; \prod_{\ell \neq p} T_{\ell}\bigl(k(\bar{x})\bigr)^{C} ,
\]
où $C$ est l'ensemble des composantes irréductibles de $\bar{D}$ et $p$ la
caractéristique résiduelle en $x$. \emph{Un facteur par branche du diviseur} :
c'est la structure entière de l'inertie modérée.

Le même calcul montre que le facteur correspondant à une composante
$\bar{D}_{c}$ coïncide avec le groupe d'inertie en n'importe quel point de
$\bar{X}'$ au-dessus du point maximal de $\bar{D}_{c}$ qui générise $\xi$ ---
et de tels points existent toujours. Donc \emph{la classe de conjugaison de ce
facteur ne dépend que de la composante irréductible de $D = X - U$}. C'est
l'énoncé qui manquait à Murre, et c'est celui que la page 7 annonçait.

\subsection*{4. La matrice des multiplicités}

Reprenons la situation de 2 avec, des deux côtés, la structure de 3 :
$\bar{V} = \bar{Y} - \operatorname{supp}\bar{E}$, $\bar{E}$ à croisements
normaux, $C'$ l'ensemble de ses composantes. Pour tout $c' \in C'$, l'image
inverse de $\bar{E}_{c'}$ est un diviseur sur $\bar{X}$ de support contenu dans
celui de $\bar{D}$, donc
\[
  \bar{f}^{*}(\bar{E}_{c'}) \;=\; \sum_{c \in C} n_{c,c'}\,\bar{D}_{c} ,
\]
les $n_{c,c'}$ étant des entiers $\geqslant 0$ non tous nuls pour $c'$ fixé.
Un calcul d'images inverses de revêtements kummériens donne alors la
commutativité de

\begin{tikzcd}[column sep=small, row sep=small, nodes={font=\scriptsize}]
G_{\xi} \arrow[r, "\varphi"] & H_{y} \\
\prod_{\ell} \underline{Z}_{\ell}(1)(\bar{x})^{C} \arrow[u] \arrow[r, "N"] &
\prod_{\ell} \underline{Z}_{\ell}(1)(\bar{y})^{C'} \arrow[u]
\end{tikzcd}

les flèches \emph{verticales} étant les épimorphismes canoniques de (3.1) et
$N$ l'homomorphisme de matrice $(n_{c,c'})$, compte tenu que
$k(\bar{y}) \to k(\bar{x})$ induit
$\underline{Z}_{\ell}(1)(\bar{y}) \xrightarrow{\ \sim\ }
\underline{Z}_{\ell}(1)(\bar{x})$.\footnote{Le tapuscrit écrit d'abord
« horizontales », corrigé de sa main en « verticales » : ce sont bien les
projections $I \to G_{\xi}$ et $I' \to H_{y}$ qui sont les épimorphismes
canoniques, la flèche horizontale du bas étant $N$ et celle du haut $\varphi$.}

\textbf{Critère.} Pour que $N$ soit un isomorphisme, il faut et il suffit que
$\operatorname{card} C = \operatorname{card} C'$ et que l'entier ordinaire
$\det N$ soit égal à $\pm\, p^{r}$, $p$ étant l'exposant caractéristique de
$k(x)$ et de $k(y)$. Dans ce cas les sous-groupes d'inertie de $H$ en $y$ sont
exactement les conjugués des images de ceux de $G$ en $x$.

Le critère est joli et il est exact : $N$ doit être inversible sur
$\prod_{\ell \neq p}\mathbb{Z}_{\ell}$, ce qui revient à demander que $\det N$
soit inversible dans cet anneau, c'est-à-dire premier à tout $\ell \neq p$,
c'est-à-dire une puissance de $p$ au signe près. La caractéristique n'entre
donc qu'à un seul endroit, et sous cette forme.

\subsection*{5. Le cas d'un diviseur image inverse}

Le numéro 5 particularise à $V = Y - \operatorname{supp} E$ avec $E$ un
diviseur $\geqslant 0$ sur $Y$ et $U = f^{-1}(V)$, sous deux conditions : que
les anneaux locaux soient noethériens aux points concernés, et que
$D = f^{*}(E)$ soit défini et régulier. Le tapuscrit s'interrompt au milieu de
cette seconde condition ; il reprend à la page 121, dans le lot 7, avec la
condition de transversalité et la conclusion sur les classes de conjugaison
d'homomorphismes $\underline{Z}_{\ell}(1) \to G$.

\end{document}
